\section*{Sets and Indices}

No nontrivial index sets are used in the implemented model. The transportation modes are:
\[
\{\text{horse},\text{bicycle},\text{handcart}\}.
\]

\section*{Parameters}

All parameter values are read from \texttt{general\_parameters.csv}.

\[
P^{\min} \; (\texttt{min\_total\_produce})
\]
minimum total produce that must be transported.

\[
\bar E \; (\texttt{max\_total\_pollution})
\]
maximum allowable total pollution.

\[
L^{\text{horse}} \; (\texttt{min\_horse\_trips})
\]
minimum required number of horse trips.

Capacities per trip:
\[
q_{\text{horse}} \; (\texttt{horse\_capacity\_per\_trip}),\qquad
q_{\text{bike}} \; (\texttt{bicycle\_capacity\_per\_trip}),\qquad
q_{\text{cart}} \; (\texttt{handcart\_capacity\_per\_trip}).
\]

Variable transportation costs per trip:
\[
c_{\text{horse}} \; (\texttt{horse\_cost\_per\_trip}),\qquad
c_{\text{bike}} \; (\texttt{bicycle\_cost\_per\_trip}),\qquad
c_{\text{cart}} \; (\texttt{handcart\_cost\_per\_trip}).
\]

Fixed usage costs:
\[
f_{\text{horse}} \; (\texttt{horse\_fixed\_cost}),\qquad
f_{\text{bike}} \; (\texttt{bicycle\_fixed\_cost}),\qquad
f_{\text{cart}} \; (\texttt{handcart\_fixed\_cost}).
\]

Pollution generated per trip:
\[
e_{\text{horse}} \; (\texttt{horse\_pollution\_per\_trip}),\qquad
e_{\text{bike}} \; (\texttt{bicycle\_pollution\_per\_trip}),\qquad
e_{\text{cart}} \; (\texttt{handcart\_pollution\_per\_trip}).
\]

Pollution charge and permit parameters:
\[
\gamma \; (\texttt{pollution\_penalty\_per\_unit}),
\qquad
\Theta \; (\texttt{horse\_pollution\_permit\_threshold}),
\qquad
F^{\text{permit}} \; (\texttt{horse\_pollution\_permit\_fee}).
\]

Big-\(M\) constant used in the implementation:
\[
M = 10000.
\]

For convenience, define total pollution
\[
E = e_{\text{horse}}x_{\text{horse}} + e_{\text{bike}}x_{\text{bike}} + e_{\text{cart}}x_{\text{cart}}.
\]

\section*{Decision Variables}

Trip counts:
\[
x_{\text{horse}} \; (\texttt{x\_horse}),\quad
x_{\text{bike}} \; (\texttt{x\_bike}),\quad
x_{\text{cart}} \; (\texttt{x\_cart}) \in \mathbb{Z}_{\ge 0}.
\]

Binary usage indicators:
\[
z_{\text{horse}} \; (\texttt{z\_horse}),\quad
z_{\text{bike}} \; (\texttt{z\_bike}),\quad
z_{\text{cart}} \; (\texttt{z\_cart}) \in \{0,1\}.
\]

Binary pollution-permit indicator:
\[
y^{\text{permit}} \; (\texttt{pollution\_permit}) \in \{0,1\}.
\]

\section*{Objective}

Minimize total transportation, fixed, pollution, and permit costs:
\[
\min \;
c_{\text{horse}}x_{\text{horse}}
+ c_{\text{bike}}x_{\text{bike}}
+ c_{\text{cart}}x_{\text{cart}}
+ f_{\text{horse}}z_{\text{horse}}
+ f_{\text{bike}}z_{\text{bike}}
+ f_{\text{cart}}z_{\text{cart}}
+ \gamma E
+ F^{\text{permit}} y^{\text{permit}}.
\]

\section*{Constraints}

Produce requirement:
\[
q_{\text{horse}}x_{\text{horse}}
+ q_{\text{bike}}x_{\text{bike}}
+ q_{\text{cart}}x_{\text{cart}}
\;\ge\;
P^{\min}.
\]

Maximum pollution:
\[
E \le \bar E.
\]

Permit-gate constraint as implemented:
\[
E \le \Theta + M\, y^{\text{permit}}.
\]

Minimum horse trips:
\[
x_{\text{horse}} \ge L^{\text{horse}}.
\]

Linking trip counts to fixed-cost binaries:
\[
x_{\text{horse}} \le M z_{\text{horse}},
\qquad
x_{\text{bike}} \le M z_{\text{bike}},
\qquad
x_{\text{cart}} \le M z_{\text{cart}}.
\]

At most one of bicycle and handcart may be used:
\[
z_{\text{bike}} + z_{\text{cart}} \le 1.
\]

Integrality and binary restrictions:
\[
x_{\text{horse}},x_{\text{bike}},x_{\text{cart}} \in \mathbb{Z}_{\ge 0},
\qquad
z_{\text{horse}},z_{\text{bike}},z_{\text{cart}},y^{\text{permit}} \in \{0,1\}.
\]

\section*{Notes on Implementation}

The formulation above matches \texttt{current\_heuristic.py} exactly.

The code introduces a binary \texttt{pollution\_permit} and charges the one-time fee \(F^{\text{permit}}\) in the objective, with the single gating constraint
\[
E \le \Theta + M\, y^{\text{permit}}.
\]
Thus, if total pollution exceeds \(\Theta\), the model must set \(y^{\text{permit}}=1\) and pay the fee. If \(E \le \Theta\), the model may set \(y^{\text{permit}}=0\), avoiding the fee. There is no reverse constraint forcing \(y^{\text{permit}}=1\) only when \(E>\Theta\); however, because the permit fee is nonnegative and minimized, the optimal solution will not set \(y^{\text{permit}}=1\) unnecessarily.

The parameter \texttt{produce\_to\_transport} is loaded by the code into \texttt{produce} but is not used in the optimization model; the implemented requirement is based on \texttt{min\_total\_produce} only.

The model is solved as a mixed-integer linear program using \texttt{GUROBI\_CMD(msg=0)}.
